(1) 圆盘半径\(r\)处速度\(\boldsymbol v=\omega r\hat\phi\)，故\(\boldsymbol v\times\boldsymbol B=\omega rB\hat r\)。沿半径积分动生电动势\(\mathcal E=\int_0^R\omega Br\,dr=k\omega\)。将外电感的反电动势一并保留，Kirchhoff方程为\(L\dot I+\mathcal RI=k\omega\)，而不是在快速变化时立即令\(I=k\omega/\mathcal R\)。
(2) 半径\(r\)处径向电流元在磁场中受切向Lorentz力，圆盘所受反向磁转矩大小为\(\int_0^R IBr\,dr=kI\)。所以\(J\dot\omega=-kI\)。两个一阶方程\(\dot\omega=-kI/J\)与\(\dot I=(k\omega-\mathcal RI)/L\)相耦合，电磁转矩自动满足能量回馈的互易系数\(k\)。
(3) 对机械方程求导并用电路方程消去\(I=-J\dot\omega/k\)，得到\(\ddot\omega+(\mathcal R/L)\dot\omega+(k^2/(JL))\omega=0\)。指数模\(e^{st}\)的特征根\(s_\pm=[-\mathcal R\pm\sqrt{\mathcal R^2-4Lk^2/J}]/(2L)\)。当\(\mathcal R^2>4Lk^2/J\)时为两个负实根；当小于时为衰减振荡；等号时出现重根和\(t e^{st}\)项，均对应电磁阻尼。
(4) 定义总储能\(E=\tfrac12J\omega^2+\tfrac12LI^2\)，求导\(J\omega\dot\omega+LI\dot I=-k\omega I+I(k\omega-\mathcal RI)=-\mathcal RI^2\le0\)。在\(L\to0\)、电路快速松弛后\(I\simeq k\omega/\mathcal R\)，于是\(J\dot\omega=-k^2\omega/\mathcal R\)，\(\omega(t)=\omega_0\exp[-k^2t/(J\mathcal R)]\)。剩余机械动能最终转成焦耳热。
